% !TeX root = ../thesis.tex
\chapter{相关技术与理论基础}

本章围绕本文后续模型库构建、统一服务封装、部署优化与原型系统验证所依赖的关键技术展开论述。与偏算法创新型论文不同，本文的研究重点不在于提出全新的知识抽取网络结构，而在于将已有代表模型组织为可训练、可封装、可部署、可监控、可比较和可选型的工程化模型库。因此，本章不仅需要说明知识抽取任务本身的形式化定义，还需要进一步阐明预训练语言模型、联合抽取范式、模型制品化、统一服务化和部署 profiling 等工程基础如何共同支撑后续章节的系统实现。

\section{知识抽取任务基础}
\subsection{任务定义与结构化表示}
面向装备制造业文本的知识抽取任务，本质上是从输入文本中识别实体、关系及其组合结构，并将其转化为能够进入知识图谱或企业知识库的结构化表示。设输入文本序列为
\begin{equation}
\mathbf{x}=\{x_{1},x_{2},\ldots,x_{n}\},
\end{equation}
其中 $x_i$ 表示第 $i$ 个 token。实体识别任务可表示为从文本序列中抽取实体集合
\begin{equation}
\mathcal{E}=\{(s_i,e_i,t_i)\mid 1\leq s_i\leq e_i\leq n\},
\end{equation}
其中 $s_i$ 与 $e_i$ 分别表示实体的起止位置，$t_i$ 表示实体类别。对于制造业场景而言，实体类别往往包括设备、部件、工艺参数、故障现象、工艺步骤、标准条款和检测指标等。

在实体识别基础上，关系抽取任务可进一步表示为
\begin{equation}
\mathcal{R}=\{(e_h,r,e_t)\mid e_h,e_t\in\mathcal{E},r\in\mathcal{L}_{r}\},
\end{equation}
其中 $e_h$ 和 $e_t$ 分别为头实体与尾实体，$r$ 表示关系类型集合 $\mathcal{L}_{r}$ 中的某一关系标签。若将联合抽取目标表示为三元组集合，则有
\begin{equation}
\mathcal{T}=\{(e_h,r,e_t)\mid e_h,e_t\in\mathcal{E}, r\in\mathcal{L}_{r}\}.
\end{equation}
该表示方式适用于 CasRel、SpERT、TPLinker、PURE、PRGC、OneRel、CopyMTL、DyGIE++ 和 GraphRel 等不同范式模型的结果对齐，是后续统一模型输出 schema 的理论基础。

从企业应用角度看，模型不仅需要输出三元组，还需要提供能够支持人工核验、结果排序和服务路由的置信度信息。基于此，本文将统一输出结果表示为
\begin{equation}
\hat{\mathcal{T}}=\{(e_h,r,e_t,p,\delta)\},
\end{equation}
其中 $p\in[0,1]$ 表示抽取结果的置信度，$\delta$ 表示证据片段或来源位置信息。该设计使抽取结果不再只是离线实验中的标签预测，而能够进一步支撑后续的结果追踪、错误分析和平台侧回溯。

\subsection{装备制造业文本特征与抽取难点}
与新闻、百科或开放域问答文本相比，装备制造业文本具有更强的领域性和更复杂的结构性。首先，制造业文本中的术语密度高，存在大量专业缩写、型号编码、工艺代号与复合参数表达，如设备部件名称、故障代码和检测条件常以混合字母数字形式出现。这一特点会直接影响分词效果与词表覆盖率，使得通用英文实现无法直接迁移到中文行业语料中。

其次，制造业文本中的关系表达具有显著的隐式性和跨句性。工艺文件往往采用流程描述或约束表达方式，同一关系可能并不在单一短句中完整出现，而需要结合上下文推断；故障报告常包含“现象—原因—处理措施”的链式描述；标准规范文本则具有严格的层级结构和长距离依赖。由此带来的问题是，单纯依赖局部上下文或短序列假设的模型在复杂文本中容易出现边界识别错误、关系遗漏和跨实体错连。

再次，制造业知识抽取通常存在关系重叠与实体重叠问题。一个设备实体可能同时参与多个关系，一个工艺参数也可能既作为条件实体又作为结果实体参与不同三元组。这使得传统“先识别实体、再分类关系”的流水线方法容易在中间阶段累积误差，进而影响最终三元组质量。因此，本文关注的模型多采用联合建模或结构化建模策略，以提升对重叠关系场景的适应能力\cite{bekoulis2018,casrel2020,tplinker2020,onerel2022}。

最后，企业平台对知识抽取模型提出的要求不止是抽取准确，还包括可复用、可部署、可回归和可治理。对于工程项目而言，若模型只能在特定实验脚本下运行、无法统一导出、无法快速比较不同配置，那么即使单次 F1 指标较好，也难以形成可长期维护的能力模块。正因为存在这一现实背景，本文在任务定义阶段就将“统一表示、统一接口和统一评测”作为基本前提，而不仅仅把知识抽取看作单次预测问题。

\subsection{任务评价方式与工程目标}
知识抽取任务的基础评价通常采用 Precision、Recall 和 F1。若预测三元组集合为 $\hat{\mathcal{T}}$，真实三元组集合为 $\mathcal{T}$，则精确率、召回率与 F1 可分别表示为
\begin{equation}
P=\frac{|\hat{\mathcal{T}}\cap\mathcal{T}|}{|\hat{\mathcal{T}}|},\qquad
R=\frac{|\hat{\mathcal{T}}\cap\mathcal{T}|}{|\mathcal{T}|},\qquad
F_{1}=\frac{2PR}{P+R}.
\end{equation}
上述指标能够衡量模型在离线标注语料上的抽取效果，是模型库纳入、中文化适配和双语实验对比的基本依据。

然而，对工程应用而言，仅以 F1 作为唯一标准是不充分的。企业知识平台面对的往往是在线接口调用、离线批量补录和资源受限部署等多类场景，模型的“优”不再等同于单一精度最优，而是要在效果、时延、资源和稳定性之间进行折中。为描述这种工程目标，本文引入面向场景的效用函数
\begin{equation}
U(m,\mathbf{p})=\alpha F_{1}(m)+\beta T(m,\mathbf{p})-\gamma L_{95}(m,\mathbf{p})-\delta M(m,\mathbf{p}),
\label{eq:foundation_utility}
\end{equation}
其中 $m$ 表示模型，$\mathbf{p}$ 表示部署配置，$T$ 表示吞吐量，$L_{95}$ 表示 P95 延迟，$M$ 表示资源开销，$\alpha,\beta,\gamma,\delta$ 为与业务场景相关的权重。式~(\ref{eq:foundation_utility}) 说明：模型选型不是孤立的算法问题，而是场景约束下的综合优化问题。

因此，本文在后续章节中将同时保留离线抽取质量指标和在线部署性能指标。一方面，第三章的模型库构建依赖离线精度与中文化适配结果；另一方面，第四章的部署优化和第五章的系统原型验证则进一步强调时延、吞吐、显存、冷启动时间和稳定性等工程指标。第二章在此先给出上述评价思想，是为了使后文不同层次的实验与系统设计共享同一理论口径。

\section{预训练模型与联合抽取方法}
\subsection{Transformer 与预训练语言模型基础}
Transformer 架构以自注意力机制为核心，能够在统一编码空间中建模长距离依赖关系\cite{vaswani2017}。相比传统的 CNN 或 BiLSTM 编码器，Transformer 更适合处理制造业文本中常见的跨短语依赖、条件约束和并列结构。BERT 通过大规模语料预训练，进一步将上下文语义表征能力引入下游任务\cite{bert2019}，从而使实体识别、关系抽取和联合抽取都能够建立在统一的预训练表示之上。

从工程实现看，预训练语言模型的重要意义并不仅在于提升精度，还在于为多模型统一适配提供共同基础。本文所纳入的多个模型虽然原始结构不同，但大量模型都依赖 BERT 类编码器进行上下文表示。这使得“统一替换主干编码器”“统一中文分词方案”“统一输入张量组织方式”成为可能，也正是后续将多模型迁移到 bert-base-chinese 并开展双语实验的前提。

在预训练语言模型发展过程中，ALBERT、SpanBERT、ERNIE 等方法分别从参数共享、span 建模和知识增强等角度改进了 BERT\cite{albert2020,spanbert2020,ernie2019}。同时，为适应资源受限环境，TinyBERT 和 MobileBERT 等轻量化模型又从蒸馏与压缩角度降低了部署成本\cite{tinybert2020,mobilebert2020}。这些工作说明，预训练模型既是抽取效果的重要来源，也是后续部署优化与模型轻量化的重要切入点。

对于本文的研究边界而言，第二章不进一步展开大语言模型生成式抽取，而是聚焦于 BERT 系列和经典联合抽取模型。这一选择具有工程合理性：一方面，项目阶段的实际目标是构建稳定、可控的知识抽取模型库；另一方面，统一制品化、统一服务接口和配置驱动 profiling 更容易在结构化抽取模型上形成可复现的工程链路。

\subsection{联合抽取模型的典型范式}
知识抽取模型从工程视角大致可以分为级联关系抽取、span-based 联合抽取、token-pair 链接建模、分阶段关系解码和图结构联合抽取等几类范式。不同范式并不只是算法上的差别，也直接决定了模型在输入输出组织、推理复杂度和部署成本上的差异。

级联关系抽取以 CasRel 为代表，其核心思想是将关系抽取转化为条件实体标注问题，即先围绕头实体识别候选关系，再进一步定位尾实体\cite{casrel2020}。该类方法链路清晰、输出结构直接、服务化封装相对容易，因此通常适合作为低时延场景中的候选模型。其工程优势在于推理逻辑可解释、错误定位较直接，但对复杂重叠关系和多跳依赖的处理仍存在局限。

Span-based 联合抽取以 SpERT 和 PURE 为代表\cite{spert2019,pure2021}。这类方法先生成或筛选候选 span，再在 span 层面进行实体分类和关系判断。与级联方法相比，span-based 方法在处理实体边界不规则、嵌套结构和复杂结果展示时更有优势，也更容易与后续知识库字段映射对齐。然而，其候选 span 生成会带来较高的组合复杂度，对长文本和较大 batch 的显存占用更敏感。

Token-pair 或 linking 型方法以 TPLinker 和 OneRel 为代表\cite{tplinker2020,onerel2022}。此类方法通过 token 对之间的关联建模来同时表达实体边界与关系链接，对实体重叠和关系重叠场景有较好适应性。其优点在于联合建模能力较强、结构表达统一，但在序列长度增加时，token 对矩阵往往带来更高的时间与空间开销，因此工程部署时需要重点关注显存峰值和吞吐下降问题。

分阶段关系解码方法以 PRGC 和 CopyMTL 为代表\cite{prgc2021,copymtl2020}。这类方法试图在任务分解与联合学习之间取得平衡，通过潜在关系筛选、多任务学习或复制机制提高三元组生成质量。相较于完全流水线方法，它们能在一定程度上降低误差传播；相较于单步联合方法，又保留了较清晰的中间状态，便于调试和服务封装。因此，该类方法非常适合在“精度与部署成本并重”的企业场景中作为折中候选。

图结构联合抽取方法以 DyGIE++ 和 GraphRel 为代表\cite{luan2018,dygiepp2019,graphrel2019}。此类方法更强调实体、关系和上下文结构之间的全局依赖建模，对长文本、多实体交互和复杂上下文约束具有较强适应能力。其工程短板在于图结构构建、候选组合和推理过程通常更复杂，导致实现与部署门槛更高，更适合作为离线补录、高精度验证或复杂场景分析模型。

为便于后续模型库构建与选型，本文将典型范式及其工程特征总结如表~\ref{tab:foundation_models_compact} 所示。
\begin{table}[htbp]
  \centering
  \caption{本文关注的知识抽取模型范式及工程特征}
  \vspace{0.5em}\wuhao
  \begin{tabularx}{\textwidth}{l l X}
    \toprule
    模型范式 & 代表模型 & 工程特征 \\
    \midrule
    级联式关系抽取 & CasRel & 推理链路清晰，接口封装直接，适合低时延在线调用场景 \\
    Span-based 联合抽取 & SpERT, PURE & 结构表达规范，便于结果展示与字段映射，但候选 span 较多时显存开销较大 \\
    Token-pair / linking & TPLinker, OneRel & 对重叠关系适应性强，统一建模能力好，但长序列下时间与空间复杂度上升明显 \\
    分阶段关系解码 & PRGC, CopyMTL & 兼顾结构信息与任务分解，便于调试与扩展，适合精度与部署成本平衡场景 \\
    图结构联合抽取 & DyGIE++, GraphRel & 全局建模能力强，适合复杂上下文和离线补录，但实现与部署复杂度相对更高 \\
    \bottomrule
  \end{tabularx}
  \label{tab:foundation_models_compact}
\end{table}

\subsection{中文化适配与双语迁移基础}
由于本文模型库中的多个模型最初面向英文数据集实现，因此在制造业中文文本场景下直接复用会面临明显的迁移问题。第一，英文实现往往默认采用英文分词器、词表和大小写处理规则，而中文文本缺少空格分隔，且术语中存在大量专业缩写、数字和符号混合表达，若不重构分词与输入处理流程，模型编码结果将与原始实现假设不一致。第二，制造业知识抽取任务中的标签体系与公开英文数据集并不完全相同，关系集合、实体边界和输出 schema 都需要重新映射。

从工程视角看，中文化适配至少涉及编码器替换、分词重构、样本生成、标签映射和训练脚本迁移五个环节。设原始模型为 $m$，中文语料集合为 $\mathcal{D}_{zh}$，中文词表与分词器为 $\mathcal{V}_{zh}$，标签体系为 $\mathcal{L}_{zh}$，则中文化适配可表示为
\begin{equation}
m^{zh}=\Phi(m,\mathcal{D}_{zh},\mathcal{V}_{zh},\mathcal{L}_{zh}),
\end{equation}
其中 $\Phi(\cdot)$ 并不只是模型权重替换操作，而是围绕数据、编码、标签和训练流程的系统性重构。

在多模型统一适配过程中，本文选择以 bert-base-chinese 作为可迁移模型的统一中文主干，一方面是因为其与 BERT 架构保持较强兼容性，便于复用原始模型结构；另一方面是因为统一主干能够降低模型库维护复杂度，使不同模型之间的输入 schema、权重导出流程和部署依赖更加接近。对于英文版本模型，则保留原始英文编码器配置，从而形成“语言不同、制品描述一致”的双语实验组织方式。这一组织思路既有利于横向对比，也为后续模型库持续扩展提供了统一接口。

\section{模型制品化与服务化基础}
\subsection{模型制品组织基础}
工程化模型库与传统实验仓库的根本差异，在于前者需要把模型从“训练结果”转变为“可交付资产”。对本文而言，一个真正可进入平台流程的模型，不应只包含权重文件，还必须同时具备元数据、标签定义、输入输出 schema、默认部署参数和版本信息。基于这一认识，本文将模型制品抽象为
\begin{equation}
\mathcal{A}=\langle w,m,s,p,v \rangle,
\label{eq:artifact_foundation}
\end{equation}
其中 $w$ 表示 ONNX 权重文件，$m$ 表示模型元数据，$s$ 表示标签与 schema 描述，$p$ 表示默认部署 Profile，$v$ 表示版本信息。只有上述对象同时完备时，模型才具备被长期管理、比较、发布和回滚的条件。

式~(\ref{eq:artifact_foundation}) 中的元数据对象尤其关键。它不仅记录模型名称、任务类型、语言版本、主干编码器、输入长度限制和制品路径，还承担“模型资产登记”的作用。对于企业平台而言，模型元数据是后续部署创建、节点调度、日志追踪和问题回溯的基础索引；对于论文原型而言，它也是实现“模型库而非模型集合”的核心组织机制。

进一步地，模型制品化也是统一实验口径的重要前提。若不同模型仍保持各自的目录结构、标签命名和输出格式，那么即便都完成了训练，也很难在同一部署系统中进行横向对比。通过将制品组织为统一对象，本文后续能够在不依赖原始训练代码细节的前提下，完成模型上传、服务启动、Profile 绑定和结果导出。这种“实验实现与工程交付解耦”的组织方式，有利于后续部署、比较与回归。

\subsection{ONNX 导出与统一推理服务基础}
ONNX 作为跨框架模型交换格式，为多模型统一部署提供了重要基础。对于本文的模型库原型而言，ONNX 的意义并不止于“能够导出模型”，而在于它使训练框架与推理框架之间形成相对稳定的接口边界。通过统一导出到 ONNX，模型可以进一步接入 ONNX Runtime、TensorRT 或其他推理后端，从而减少因原始实现环境差异带来的部署阻碍。

在服务层，统一推理接口是屏蔽模型内部差异的关键。本文将服务请求抽象为
\begin{equation}
q=\langle \texttt{text},\texttt{scene},\texttt{model\_id},\texttt{profile}\rangle,
\label{eq:service_request_foundation}
\end{equation}
将服务输出抽象为
\begin{equation}
y=\langle \mathcal{E},\mathcal{R},\hat{\mathcal{T}},\texttt{latency},\texttt{version}\rangle.
\label{eq:service_response_foundation}
\end{equation}
式~(\ref{eq:service_request_foundation}) 和式~(\ref{eq:service_response_foundation}) 表明，平台侧只需要理解统一接口，而不必感知模型内部究竟采用了级联解码、span 分类还是 token-pair 标注。这种封装方式使模型切换、版本回归和场景路由都可以在服务层完成，而不必反复修改调用逻辑。

若进一步考虑证据追踪与人工核验需求，统一服务接口还可以把模型推理时间、版本号、节点信息和部署 Profile 作为结果附加字段返回。这样做的直接收益在于：系统不仅能回答“抽取了什么”，还能够回答“由哪个模型、在何种配置下抽取出来”。对于制造业知识图谱建设而言，这种可追踪性本身就是平台治理的一部分。

\begin{figure}[htbp]
  \centering
  \begin{tikzpicture}[
    box/.style={draw, rectangle, text width=2.35cm, minimum height=0.95cm, align=center, inner sep=3pt, font=\wuhao},
    arr/.style={-Latex, semithick}
  ]
    \node[box] (a) at (0,0) {训练模型};
    \node[box] (b) at (3.0,0) {ONNX 制品};
    \node[box] (c) at (6.0,0) {统一推理服务};
    \node[box] (d) at (9.0,0) {部署与 profiling};
    \node[box] (e) at (12.0,0) {平台调用与回溯};
    \draw[arr] (a) -- (b);
    \draw[arr] (b) -- (c);
    \draw[arr] (c) -- (d);
    \draw[arr] (d) -- (e);
  \end{tikzpicture}
  \caption{知识抽取模型从训练结果到平台调用的工程链路}
  \label{fig:foundation_service_chain}
\end{figure}

\subsection{模型库治理与版本化管理基础}
模型服务真正进入工程流程后，仅有“能运行”仍然不够，还必须具备持续治理能力。这里的治理主要包括三个方面：其一，版本管理，即记录模型从训练权重、ONNX 导出到部署实例的完整版本链路；其二，发布回归，即在新版本上线前比较其精度、延迟、显存和稳定性是否满足门槛；其三，问题回溯，即能够根据日志、指标和任务历史追踪到具体模型版本和部署配置。

对于多模型并存的模型库而言，治理的重要性甚至高于单模型最优。因为企业平台真正面临的问题不是“某个模型是否最好”，而是“当业务场景、算力约束和数据来源变化时，应如何在多个候选模型之间完成比较、选择和切换”。这就要求模型库在设计之初就保留统一元数据字段、统一部署模板和统一结果记录方式，从而为后续的选型机制提供可计算、可回溯的依据。

因此，本文后续章节中的发布回归、Profile 绑定、日志监控与结果导出，并不是附加功能，而是模型库治理能力的具体表现。第二章在此给出这一基础，是为了说明为什么本文研究对象必须是“模型库原型系统”，而不能仅停留在“若干模型实验结果”的层面。

\section{部署评测与资源约束基础}
\subsection{面向企业应用的性能指标体系}
对于企业知识平台中的知识抽取服务，模型质量不能只由离线 F1 决定，还必须同时考虑延迟、吞吐量、显存占用、冷启动时间和稳定性等指标。部署评测指标可分为抽取质量、服务性能、资源开销、稳定性和治理能力五类，如表~\ref{tab:foundation_metrics_compact} 所示。

\begin{table}[htbp]
  \centering
  \caption{面向工程部署的核心评测指标}
  \vspace{0.5em}\wuhao
  \begin{tabularx}{\textwidth}{l X}
    \toprule
    指标类别 & 指标含义 \\
    \midrule
    抽取质量 & Precision、Recall、F1，用于评价模型输出结果对业务知识构建的有效性 \\
    服务性能 & 平均延迟、P95 延迟、P99 延迟、吞吐量，用于评价在线服务与批处理能力 \\
    资源开销 & GPU 显存峰值、GPU 利用率、CPU 内存和 CPU 利用率，用于衡量部署成本 \\
    稳定性 & 成功率、超时率、异常重启次数、冷启动时间，用于评价长期运行能力 \\
    治理能力 & 版本追踪、日志完整性、结果可回溯性，用于评价模型服务的可维护性 \\
    \bottomrule
  \end{tabularx}
  \label{tab:foundation_metrics_compact}
\end{table}

其中，平均延迟主要反映系统的一般响应速度，而 P95、P99 等尾延迟指标更能反映服务在高并发或波动条件下的最坏体验。对于实时故障抽取场景，尾延迟往往比平均延迟更有意义；对于离线批量知识补录场景，吞吐量和吞吐/显存比则更值得关注。冷启动时间则主要影响新实例发布和扩容时的服务可用性，是工程部署中经常被忽视但实际非常关键的指标。

此外，稳定性与治理能力指标虽然不直接提升抽取效果，却决定了模型能否长期进入平台运行。若某模型在压力测试下成功率不稳定，或部署后缺乏版本追踪和日志记录能力，那么其工程价值将明显受限。因此，本文在第四章和第五章中会将性能指标与治理能力一并纳入比较框架，而不是仅报告单组 benchmark 数值。

\subsection{推理后端与部署基础设施}
在模型统一导出为 ONNX 之后，推理后端的选择将直接影响系统性能与部署复杂度。基于 PyTorch 的原始推理环境具有实现直接、调试方便等优点，但在服务化部署中往往面临依赖环境重、冷启动时间长和跨模型统一管理困难等问题。围绕模型服务系统的已有研究表明，推理运行时、批处理策略、资源配置和尾延迟控制是影响部署效果的关键因素\cite{clockwork2020,clipper2017}。基于这一认识，本文在原型实现中选用 ONNX Runtime 作为核心推理运行时，并结合图优化、线程配置、量化和多后端执行等机制组织统一推理链路。

当部署进一步面向 GPU 优化时，精度压缩、算子融合和动态输入配置通常能够改善推理效率，但也会引入引擎生成、输入范围设定和缓存管理等额外复杂度。因此，在实际系统设计中，需要根据模型类别、输入长度和场景需求决定是否启用更激进的 GPU 优化策略。

对于服务编排与可观测性而言，现有模型服务系统研究普遍将 profiling、资源观测和服务等级约束视为部署优化的重要组成部分\cite{clipper2017,clockwork2020}。本文虽然未将整套通用推理服务器完整纳入原型系统，但在部署 profiling 思想上借鉴了“通过配置文件扫描 batch、并发、后端和实例配置”的方式；在系统实现中，则结合 GPUStack 完成模型部署管理，并通过日志、资源指标和基准测试结果支撑部署对比。上述基础设施共同说明：部署优化不是单一模型参数调节问题，而是模型、运行时和服务系统协同作用的结果。

\subsection{Profile 驱动的性能 profiling 思想}
在工程实践中，所谓“写几个 profile 文件”并不应被理解为随意填写几组部署参数，而应理解为用配置文件显式描述一组可重复执行的 profiling 方案。设部署 Profile 表示为
\begin{equation}
\mathbf{p}=[b,c,l,e,r,\tau],
\label{eq:foundation_profile}
\end{equation}
其中 $b$ 表示 batch size，$c$ 表示并发数，$l$ 表示最大输入长度，$e$ 表示精度模式，$r$ 表示推理后端，$\tau$ 表示线程或执行环境配置。任何一组 Profile 都需要满足资源约束：
\begin{equation}
M_{\text{static}}+M_{\text{act}}(\mathbf{p})+M_{\text{ws}}(\mathbf{p})\leq \rho M_{\text{gpu}},
\label{eq:foundation_memory}
\end{equation}
其中 $M_{\text{static}}$ 表示模型静态显存占用，$M_{\text{act}}$ 表示激活与输入相关开销，$M_{\text{ws}}$ 表示运行时工作空间，$\rho$ 为安全冗余系数。

Profile 驱动 profiling 的意义主要体现在三个方面。第一，它将原本散落在脚本命令行中的实验环境显式化，使不同模型、不同版本和不同后端的测试条件能够被准确复现。第二，它使“延迟优先、吞吐优先、资源优先、回归验证”这几类典型部署目标能够对应到不同模板，而不是混在同一套 benchmark 中。第三，它为后续模型比较和选型提供统一输入，使模型选型过程从经验判断转变为有数据依据的配置搜索过程。

受 Clipper 与 Clockwork 等系统工作启发\cite{clipper2017,clockwork2020}，本文将部署优化理解为从模型效果到服务 SLO 的完整链路问题，而不是孤立的算子加速问题。在知识抽取模型库场景下，Profile 不只是“跑得更快”的调参文件，更是连接模型资产、部署后端、基准测试和选型规则的统一配置对象。这一认识将直接支撑后文第四章的部署优化方法和第五章的基准测试模块实现。

\section{本章小结}
本章围绕知识抽取模型库原型所依赖的任务基础、模型方法、制品化与服务化基础以及部署评测基础进行了系统梳理。首先，从知识抽取任务定义出发，说明了实体、关系和三元组的统一表示方式，以及制造业文本在术语密度、关系重叠、跨句依赖和工程约束方面的特点。其次，从 Transformer 与预训练语言模型出发，分析了级联式、span-based、token-pair、分阶段解码和图结构联合抽取等典型范式的工程差异，并说明了中文化适配与双语迁移的必要性。再次，围绕模型制品化与统一推理服务，给出了模型资产组织、ONNX 导出、统一服务接口和版本治理的基础概念。最后，从部署评测与资源约束角度，阐明了多维性能指标体系、推理后端选择以及 Profile 驱动 profiling 的理论依据。

第二章为后续章节提供统一的技术基础。第三章的模型库构建以前文的模型范式分析与中文化适配为前提；第四章的工程化部署以前文的制品化、推理后端和 Profile 理念为支撑；第五章的原型系统验证则以统一接口、统一指标和统一治理口径为依据。本章构成了全文从“模型研究”走向“系统原型验证”的理论与工程基础。
